\documentclass[12pt]{article}
\usepackage[utf8]{inputenc}
\usepackage[T1]{fontenc}
\usepackage[margin=1in]{geometry}
\usepackage{setspace}
\usepackage{xcolor}
\usepackage[colorlinks=true, urlcolor=blue!60!black, linkcolor=blue!60!black]{hyperref}
\setlength{\parindent}{2em}
\setlength{\parskip}{1.25em}
\renewcommand{\baselinestretch}{1.0}
\begin{document}
\begin{center}
{\LARGE \textit{New York v. United States}, 505 U.S. 144 (1992)\par}
\vspace{0.8em}
{\large Case Brief\par}
\end{center}
\vspace{1em}

\section{Memory Jogger}

While Congress may encourage states to regulate through conditional federal spending or conditional preemption under the Commerce Clause, the Tenth Amendment's anti-commandeering principle forbids Congress from compelling state legislatures to enact or administer a federal regulatory program, rendering the Low-Level Radioactive Waste Policy Amendments Act's ``take title'' provision unconstitutional (\href{https://www.courtlistener.com/opinion/112768/new-york-v-united-states/}{CourtListener}).

\section{Detailed Case Facts}

By 1979, after commercial low-level radioactive waste disposal sites in Illinois, Kentucky, and New York closed, only three commercial disposal sites remained open nationwide: Beatty, Nevada; Hanford, Washington; and Barnwell, South Carolina (\href{https://www.law.cornell.edu/supremecourt/text/505/144}{Cornell LII}). When temporary shutdowns in Nevada and Washington forced South Carolina to shoulder the nation's waste, South Carolina's Governor cut accepted volume by 50\%, and Washington and Nevada threatened permanent closure (\href{https://www.law.cornell.edu/supremecourt/text/505/144}{Cornell LII}). Acting largely on proposals from the National Governors' Association, Congress passed the Low-Level Radioactive Waste Policy Act of 1980 (Pub. L. 96-573), declaring each state responsible for disposing of low-level radioactive waste generated within its borders and authorizing congressionally approved regional compacts to exclude out-of-region waste beginning in 1986 (\href{https://www.law.cornell.edu/supremecourt/text/505/144}{Cornell LII}).

When 31 states still lacked approved disposal outlets by 1985, Congress enacted a state-brokered compromise: the Low-Level Radioactive Waste Policy Amendments Act of 1985, Pub. L. 99-240, 99 Stat. 1842, 42 U.S.C. \S 2021b et seq. (\href{https://www.law.cornell.edu/supremecourt/text/505/144}{Cornell LII}). The three sited states agreed to accept out-of-region waste for seven additional years (1986--1992) while charging graduated surcharges (\$10, \$20, and \$40 per cubic foot), and the Act created three sets of ``incentives'' to induce unsited states to meet statutory siting milestones (\href{https://www.law.cornell.edu/supremecourt/text/505/144}{Cornell LII}):
\begin{enumerate}
    \item \textbf{Monetary Incentives (\S 2021e(d)(2)(A)--(B)):} One quarter of the surcharges collected by sited states was placed into a Secretary of Energy escrow account and rebated to states that met statutory milestones (1986, 1988, 1990, and 1993).
    \item \textbf{Access Incentives (\S 2021e(e)(2)):} Sited states and compacts were authorized to levy escalating surcharges (double, triple, or quadruple) and eventually deny disposal-site access altogether to waste generated in noncomplying states.
    \item \textbf{The ``Take Title'' Provision (\S 2021e(d)(2)(C)):} Any state or compact region unable to provide for disposal of all internally generated waste by January 1, 1996, was required, upon request of the waste's generator or owner, to take title to and possession of the waste and assume liability for all direct and indirect damages caused by failing to take possession promptly.
\end{enumerate}

New York, a major generator of low-level radioactive waste, declined to join a regional compact and instead enacted state legislation to site an in-state disposal facility (\href{https://www.law.cornell.edu/supremecourt/text/505/144}{Cornell LII}). When state officials identified five potential sites---three in Allegany County and two in Cortland County---local residents mounted fierce political opposition (\href{https://www.law.cornell.edu/supremecourt/text/505/144}{Cornell LII}).

\section{Procedural History}

In 1990, the State of New York, Allegany County, and Cortland County filed suit against the United States in the U.S. District Court for the Northern District of New York seeking a declaratory judgment that the 1985 Act violated the Tenth Amendment, the Eleventh Amendment, the Fifth Amendment Due Process Clause, and the Article IV, \S 4 Guarantee Clause; the sited states of Washington, Nevada, and South Carolina intervened as defendants (\href{https://www.law.cornell.edu/supremecourt/text/505/144}{Cornell LII}). The District Court dismissed the complaint (\textit{New York v. United States}, 757 F. Supp. 10 (N.D.N.Y. 1990)), and the U.S. Court of Appeals for the Second Circuit affirmed (\textit{New York v. United States}, 942 F.2d 114 (2d Cir. 1991)) (\href{https://tile.loc.gov/storage-services/service/ll/usrep/usrep505/usrep505144/usrep505144.pdf}{Library of Congress U.S. Reports}). Petitioners abandoned their Fifth and Eleventh Amendment claims on appeal, pursuing only the Tenth Amendment and Guarantee Clause challenges (\href{https://www.law.cornell.edu/supremecourt/text/505/144}{Cornell LII}). The Supreme Court granted certiorari (Nos. 91-543, 91-558, and 91-563), heard oral argument on March 30, 1992, and issued its decision on June 19, 1992 (\href{https://www.oyez.org/cases/1991/91-543}{Oyez}; \href{https://www.law.cornell.edu/supremecourt/text/505/144}{Cornell LII}).

\section{Judicial Votes}

The Court ruled \textbf{9--0} to uphold the monetary and access incentives and \textbf{6--3} to invalidate and sever the ``take title'' provision (\href{https://www.law.cornell.edu/supremecourt/text/505/144}{Cornell LII}):
\begin{itemize}
    \item \textbf{Majority (6 in full; 9 as to Parts III-A and III-B):} Justice O'Connor delivered the opinion of the Court, joined in full by Chief Justice Rehnquist and Justices Scalia, Kennedy, Souter, and Thomas, and joined as to Parts III-A (monetary incentives) and III-B (access incentives) by Justices White, Blackmun, and Stevens.
    \item \textbf{Concurring in Part and Dissenting in Part (3):} Justice White filed an opinion concurring in part and dissenting in part, joined by Justices Blackmun and Stevens (joining Parts III-A and III-B, but dissenting from the invalidation of the take-title provision and the remainder of the opinion).
    \item \textbf{Concurring in Part and Dissenting in Part (1):} Justice Stevens also filed a separate opinion concurring in part and dissenting in part, joining Justice White's partial dissent and rejecting the anti-commandeering rule altogether.
\end{itemize}

\section{Holding}

\textbf{Judgment:} Affirmed in part and reversed in part (\href{https://www.law.cornell.edu/supremecourt/text/505/144}{Cornell LII}). The monetary incentives (\S 2021e(d)(2)(A)--(B)) and access incentives (\S 2021e(e)(2)) are constitutional exercises of Congress's Spending and Commerce Clause powers, but the ``take title'' provision (\S 2021e(d)(2)(C)) violates the Tenth Amendment and is severable from the remainder of the Act (\href{https://www.law.cornell.edu/supremecourt/text/505/144}{Cornell LII}).

\textbf{Rule Announced:} Congress may not ``commandeer'' the legislative processes of the states by directly compelling them to enact and enforce a federal regulatory program (\href{https://www.law.cornell.edu/supremecourt/text/505/144}{Cornell LII}). While Congress may encourage state regulation through conditional grants under the Spending Clause (satisfying \textit{South Dakota v. Dole}) or by offering states a choice between regulating private activity under federal standards or facing direct federal preemption under the Commerce Clause, Congress may not force states to choose between two unconstitutionally coercive alternatives (\href{https://www.law.cornell.edu/supremecourt/text/505/144}{Cornell LII}). Moreover, state officials' political consent to a federal statute cannot ratify an unconstitutional departure from the Constitution's federal structure (\href{https://www.law.cornell.edu/supremecourt/text/505/144}{Cornell LII}).

\textbf{Holding vs. Dicta:} The holdings are that (1) the monetary incentives are valid under the Commerce, Taxing, and Spending Clauses; (2) the access incentives are valid conditional exercises of the Commerce Clause; (3) the take-title provision is unconstitutional commandeering under the Tenth Amendment; (4) New York's prior political support does not waive or estop its Tenth Amendment claim; and (5) the take-title provision is severable (\href{https://www.law.cornell.edu/supremecourt/text/505/144}{Cornell LII}). By contrast, Justice O'Connor's extended discussion suggesting that perhaps not all Guarantee Clause claims are nonjusticiable political questions (while assuming justiciability only to reject petitioners' claim on the merits) is dicta (\href{https://www.law.cornell.edu/supremecourt/text/505/144}{Cornell LII}).

\section{Analysis of Opinions}

\subsection{Majority Opinion (Justice O'Connor)}

Justice O'Connor framed the Tenth Amendment and Article I's enumerated powers as ``mirror images of each other'': whether one asks if Congress exceeded its delegated powers or invaded state sovereignty reserved by the Tenth Amendment, the structural limit is the same (\href{https://www.law.cornell.edu/supremecourt/text/505/144}{Cornell LII}). Distinguishing \textit{Garcia v. San Antonio Metropolitan Transit Authority}, 469 U.S. 528 (1985), as involving generally applicable federal laws binding both states and private parties, she turned to whether Congress may use states ``as implements of regulation'' (\href{https://www.law.cornell.edu/supremecourt/text/505/144}{Cornell LII}). Reviewing \textit{Hodel v. Virginia Surface Mining \& Reclamation Ass'n} (1981), \textit{FERC v. Mississippi} (1982), and the Constitutional Convention's rejection of the New Jersey Plan in favor of the Virginia Plan---along with statements by Hamilton, Madison, Randolph, Ellsworth, and Pinckney---she concluded that the Framers explicitly chose a Constitution empowering Congress to regulate individuals directly rather than coercing states as political bodies (\href{https://www.law.cornell.edu/supremecourt/text/505/144}{Cornell LII}). She grounded this rule in political accountability: when Congress encourages states via conditional spending or conditional preemption, state voters retain the ultimate choice and know whom to hold responsible; when Congress compels state legislation, state officials bear local blame while federal officials remain insulated (\href{https://www.law.cornell.edu/supremecourt/text/505/144}{Cornell LII}).

Applying this framework, Justice O'Connor first construed \S 2021c(a)(1)(A) under constitutional avoidance as a set of three incentives rather than a freestanding mandate (\href{https://www.law.cornell.edu/supremecourt/text/505/144}{Cornell LII}). She upheld the \textbf{monetary incentives} (Part III-A) as valid exercises of Congress's powers to authorize state burdens on interstate commerce, impose a federal tax, and condition segregated-fund grants under \textit{South Dakota v. Dole}, 483 U.S. 203 (1987) (\href{https://www.law.cornell.edu/supremecourt/text/505/144}{Cornell LII}). She upheld the \textbf{access incentives} (Part III-B) as a valid conditional exercise of the Commerce Clause: if a state declined to regulate, any burden fell on private waste generators losing out-of-state disposal access rather than on the state as a sovereign (\href{https://www.law.cornell.edu/supremecourt/text/505/144}{Cornell LII}).

By contrast, Justice O'Connor invalidated the \textbf{take-title provision} (Part III-C), concluding that Congress ``crossed the line distinguishing encouragement from coercion'' (\href{https://www.law.cornell.edu/supremecourt/text/505/144}{Cornell LII}). Neither option offered to the states was constitutionally permissible standing alone: forcing states to take title and assume generators' damages was an unconstitutional compelled subsidy commandeering state governments, and ordering states to legislate according to Congress's instructions was an unconstitutional direct command (\href{https://www.law.cornell.edu/supremecourt/text/505/144}{Cornell LII}). Offering ``[a] choice between two unconstitutionally coercive regulatory techniques is no choice at all'' (\href{https://www.law.cornell.edu/supremecourt/text/505/144}{Cornell LII}). She rejected the federal government's arguments that a strong federal interest, state-court Supremacy Clause enforcement (\textit{Testa v. Katt}), federal judicial injunctions (\textit{Ex parte Young}), or congressional arbitration of interstate disputes could override this categorical bar (\href{https://www.law.cornell.edu/supremecourt/text/505/144}{Cornell LII}). She also rejected the sited states' consent and estoppel defenses: federalism protects individual liberty rather than state officials' personal political convenience, so state officials cannot consent to enlarge federal power any more than a President can consent to a separation-of-powers violation (\href{https://www.law.cornell.edu/supremecourt/text/505/144}{Cornell LII}). Finally, she rejected petitioners' Guarantee Clause challenge to the two valid incentives and held the take-title provision severable under \textit{Alaska Airlines, Inc. v. Brock}, 480 U.S. 678 (1987), because the remaining incentives preserved the Act's core purpose (\href{https://www.law.cornell.edu/supremecourt/text/505/144}{Cornell LII}).

\subsection{Opinion Concurring in Part and Dissenting in Part (Justice White, joined by Justices Blackmun and Stevens)}

Justice White concurred in upholding the monetary and access incentives (Parts III-A and III-B) but dissented from the invalidation of the take-title provision (\href{https://www.law.cornell.edu/supremecourt/text/505/144}{Cornell LII}). He argued that the majority stripped the 1985 Act of its historical context: the 1980 and 1985 statutes were not top-down federal impositions, but state-initiated interstate bargains negotiated through the National Governors' Association and ratified by Congress acting as an interstate referee (\href{https://www.law.cornell.edu/supremecourt/text/505/144}{Cornell LII}). Because New York actively supported the 1985 Act and reaped the benefit of seven extra years of waste exports to sited states, Justice White contended that New York's conduct manifested assent to an interstate agreement under the Compact Clause (Art. I, \S 10, cl. 3) and should estop the State from repudiating its bargain (\href{https://www.law.cornell.edu/supremecourt/text/505/144}{Cornell LII}).

Doctrinally, Justice White criticized the majority's distinction between generally applicable laws and state-directed mandates as unsupported by \textit{Garcia}, \textit{Hodel}, and \textit{FERC}, noting that \textit{Hodel}'s anti-commandeering language was dicta and that \textit{Garcia} properly left federalism protections primarily to the national political process---which functioned fully here (\href{https://www.law.cornell.edu/supremecourt/text/505/144}{Cornell LII}). He also pointed out the irony that under the Court's own logic, Congress could achieve virtually the same result through Spending Clause conditions or direct federal preemption excluding state input, thereby discouraging Congress from deferring to cooperative state-led solutions (\href{https://www.law.cornell.edu/supremecourt/text/505/144}{Cornell LII}).

\subsection{Opinion Concurring in Part and Dissenting in Part (Justice Stevens)}

Justice Stevens joined Parts III-A and III-B of the Court's opinion and Justice White's partial dissent, writing separately to reject the anti-commandeering doctrine at its root (\href{https://www.law.cornell.edu/supremecourt/text/505/144}{Cornell LII}). He argued that under the Articles of Confederation, Congress possessed the power to issue commands to the states, and the Constitution enhanced---rather than diminished---federal authority by adding power to regulate individuals directly (\href{https://www.law.cornell.edu/supremecourt/text/505/144}{Cornell LII}). Observing that the federal government routinely directs state functions (such as state-operated railroads, schools, prisons, and elections) and that the Supreme Court itself can order states to remedy interstate nuisances under interstate common law, Justice Stevens concluded that neither the Tenth Amendment nor constitutional structure forbids Congress from commanding states to implement federal waste-disposal standards (\href{https://www.law.cornell.edu/supremecourt/text/505/144}{Cornell LII}).

\section{Examples: Future Applications}

\begin{enumerate}
    \item \textbf{Same-Side (Unconstitutional Commandeering): Federal Mandate Requiring State Legislatures to Enact Carbon-Cap Statutes.} Congress enacts a statute declaring that every state legislature must pass a state law capping industrial carbon emissions by 40\% by 2030, or else the state government must pay compensatory damages to residents harmed by air pollution. Like the take-title provision in \textit{New York v. United States}, this law offers states a ``choice'' between two unconstitutionally coercive federal commands---legislating according to Congress's instructions or assuming private-party damages---and directly commandeers state legislative processes (\href{https://www.law.cornell.edu/supremecourt/text/505/144}{Cornell LII}).

    \item \textbf{Same-Side (Unconstitutional Commandeering): Federal Statute Prohibiting States from Repealing State Laws on Sports Gambling.} Congress passes a law barring state legislatures from ``authorizing'' or licensing sports betting, thereby preventing states from repealing their existing state-law prohibitions on sports gambling. As the Supreme Court later held in \textit{Murphy v. National Collegiate Athletic Ass'n}, 584 U.S. 453 (2018), applying \textit{New York v. United States}, there is no constitutional difference between compelling a state legislature to enact a law and prohibiting it from repealing one; both issue direct orders to state legislatures in violation of the anti-commandeering principle (\href{https://www.oyez.org/cases/2017/16-476}{Oyez (Murphy)}).

    \item \textbf{Opposite-Side (Constitutional under \textit{New York}): Cooperative Federalism with Conditional Preemption.} Congress enacts a Clean Air Act amendment directing that if a state declines to adopt a state implementation plan meeting federal emissions standards, the Environmental Protection Agency will directly regulate private factories within that state under a federal plan. Under \textit{New York v. United States} (Part II-B-2 and Part III-B), this is classic conditional preemption: the state retains the sovereign choice not to regulate or spend state funds, and any regulatory burden falls directly on private emitters rather than the state government (\href{https://www.law.cornell.edu/supremecourt/text/505/144}{Cornell LII}).

    \item \textbf{Opposite-Side (Constitutional under \textit{New York}): Generally Applicable Regulation of States as Market Participants.} Congress passes a federal data-privacy statute prohibiting both private data brokers and state departments of motor vehicles from selling drivers' personal information without consent. As \textit{New York v. United States} distinguished generally applicable statutes (and as the Court later confirmed in \textit{Reno v. Condon}, 528 U.S. 141 (2000)), regulating states in their capacity as database owners alongside private actors does not compel states in their sovereign capacity to regulate their own citizens (\href{https://journals.law.harvard.edu/jol/2025/02/22/anticommandeering-indian-affairs-legislation/}{Harvard Journal on Legislation}).

    \item \textbf{Fence-Sitter (Close / Contested under \textit{New York}): Federal Statute Imposing Record-Keeping and Procedural Duties on State Child-Welfare Agencies and Courts.} Congress enacts a federal statute governing state-court adoption proceedings involving a federally protected class of children, requiring state agencies to make ``active efforts'' to prevent family breakup and requiring state courts to maintain and transmit placement records to a federal agency. Challengers argue that imposing affirmative operational and record-keeping duties on state executive agencies commandeers state officials in violation of \textit{New York} and \textit{Printz v. United States}, 521 U.S. 898 (1997); defenders argue that the substantive standards apply evenhandedly to both public and private petitioners and that record-keeping duties imposed on state courts are permissible ancillary tasks under the Supremacy Clause. Because the boundary between permissible Supremacy Clause directives to state courts, evenhanded regulation, and impermissible executive commandeering is finely drawn (as litigated in \textit{Haaland v. Brackeen}, 599 U.S. 255 (2023)), such hybrid statutes sit squarely on the doctrinal fence (\href{https://journals.law.harvard.edu/jol/2025/02/22/anticommandeering-indian-affairs-legislation/}{Harvard Journal on Legislation}; \href{https://www.law.cornell.edu/supct/html/95-1478.ZO.html}{Cornell LII (Printz)}).
\end{enumerate}

\section{Critique}

\textit{New York v. United States} inaugurated the modern anti-commandeering doctrine---later extended to state executive officers in \textit{Printz v. United States}, 521 U.S. 898 (1997), and to negative prohibitions on state legislatures in \textit{Murphy v. NCAA}, 584 U.S. 453 (2018)---and drew searching critique from both originalist and progressive scholars (\href{https://www.law.cornell.edu/supct/html/95-1478.ZO.html}{Cornell LII (Printz)}; \href{https://www.oyez.org/cases/2017/16-476}{Oyez (Murphy)}). From an \textbf{originalist and textualist perspective}, defenders praised Justice O'Connor for reviving structural federalism and grounding the anti-commandeering rule in the Founding shift from the Articles of Confederation to direct national legislation over individuals (\href{https://scholarlycommons.law.northwestern.edu/cgi/viewcontent.cgi?article=1044&context=nulr_online}{Northwestern University Law Review Colloquy}). However, textualist and formalist critics noted that Justice O'Connor candidly conceded the limit ``is not derived from the text of the Tenth Amendment itself,'' which she called ``essentially a tautology,'' prompting scholars such as Martin H. Redish and others to question whether the Court derived a rigid categorical bar from structural inference and historical silence rather than clear constitutional text (\href{https://journals.law.harvard.edu/jol/2025/02/22/anticommandeering-indian-affairs-legislation/}{Harvard Journal on Legislation}; \href{https://scholarship.law.nd.edu/cgi/viewcontent.cgi?article=4931&context=ndlr}{Notre Dame Law Review}).

From a \textbf{progressive and functionalist perspective}, scholars challenged both the formalist rigidity and the empirical premises of \textit{New York}. Commentators scrutinized the Court's ``political accountability'' rationale, arguing that voters are often well aware when state officials act under federal compulsion and that cooperative federalism schemes (like conditional spending and conditional preemption, which \textit{New York} expressly preserved) can blur lines of accountability just as much as direct mandates (\href{https://scholarlycommons.law.northwestern.edu/cgi/viewcontent.cgi?article=1329&context=nulr}{Northwestern University Law Review}). Environmental and federalism scholars also echoed Justice White's practical critique: by invalidating the enforcement backstop of a state-negotiated interstate compromise, \textit{New York} incentivizes Congress either to preempt states entirely---centralizing power in Washington---or to rely on aggressive Spending Clause conditions, potentially undermining the very state autonomy federalism aims to protect (\href{https://digitalcommons.law.villanova.edu/cgi/viewcontent.cgi?article=1254&context=elj}{Villanova Environmental Law Journal}; \href{https://brooklynworks.brooklaw.edu/context/faculty/article/2006/viewcontent/Gold22.pdf}{Brooklyn Law School}).

At its core, \textit{New York v. United States} turns on \textbf{competing constitutional values}: preserving dual sovereignty, democratic accountability, and structural checks against concentrated federal power, versus enabling pragmatic national solutions to interstate collective-action problems (like radioactive waste disposal) through flexible cooperative bargains between Congress and the states.

\section{Key Quotations}

\begin{quote}
``We conclude that while Congress has substantial power under the Constitution to encourage the States to provide for the disposal of the radioactive waste generated within their borders, the Constitution does not confer upon Congress the ability simply to compel the States to do so.''
\end{quote}
\noindent --- Justice O'Connor, Opinion of the Court (\href{https://www.law.cornell.edu/supremecourt/text/505/144}{Cornell LII}).

\begin{quote}
``In providing for a stronger central government, therefore, the Framers explicitly chose a Constitution that confers upon Congress the power to regulate individuals, not States.''
\end{quote}
\noindent --- Justice O'Connor, Opinion of the Court (\href{https://www.law.cornell.edu/supremecourt/text/505/144}{Cornell LII}).

\begin{quote}
``A choice between two unconstitutionally coercive regulatory techniques is no choice at all. Either way, `the Act commandeers the legislative processes of the States by directly compelling them to enact and enforce a federal regulatory program,' an outcome that has never been understood to lie within the authority conferred upon Congress by the Constitution.''
\end{quote}
\noindent --- Justice O'Connor, Opinion of the Court (\href{https://www.law.cornell.edu/supremecourt/text/505/144}{Cornell LII}).

\begin{quote}
``States are not mere political subdivisions of the United States. State governments are neither regional offices nor administrative agencies of the Federal Government. The positions occupied by state officials appear nowhere on the Federal Government's most detailed organizational chart.''
\end{quote}
\noindent --- Justice O'Connor, Opinion of the Court (\href{https://www.law.cornell.edu/supremecourt/text/505/144}{Cornell LII}).

\begin{quote}
``The ultimate irony of the decision today is that in its formalistically rigid obeisance to `federalism,' the Court gives Congress fewer incentives to defer to the wishes of state officials in achieving local solutions to local problems.''
\end{quote}
\noindent --- Justice White, Opinion Concurring in Part and Dissenting in Part, joined by Justices Blackmun and Stevens (\href{https://www.law.cornell.edu/supremecourt/text/505/144}{Cornell LII}).

\end{document}
